\section{Generative Models in Remote Sensing}
\label{sec:generative_rs}
Generative Adversarial Networks (GANs)~\cite{Goodfellow_2014} have found significant applications in Earth Observation, particularly with SAR and coherent imaging data.
These applications provide the direct inspiration for the methodology adopted in this thesis~\cite{Li2022SRGANRemoteSensing,yebasse2025advances}.

In the context of RS data, GANs offer an attractive framework for learning complex, high-dimensional distributions of radargrams without requiring explicit parametric models of clutter, speckle, or instrumental effects~\cite{yebasse2025advances}.
Instead of hand-crafting statistical models for each noise source, a GAN can be trained to reproduce the joint distribution of real RS observations conditioned on auxiliary inputs (e.g., simulations, masks), implicitly capturing the characteristic speckle patterns, amplitude statistics, and large-scale clutter geometries.
This flexibility comes at the cost of more delicate training dynamics, which motivates the careful choice of architectures and loss functions discussed in Chapter~\ref{cha:methodology}.

\subsection{Despeckling and Super-Resolution in Coherent Imaging}
SAR images, like radargrams, suffer from multiplicative speckle noise.
GAN-based models have demonstrated superior performance in despeckling compared to classical filters (e.g., Lee or Frost filters)~\cite{Ma2020SARGAN,Gu2019GANDespeckling}.
By employing a discriminator that judges texture realism, GANs can hallucinate high-frequency details that are statistically consistent with the terrain, rather than blurring them out.
This leads to outputs that are visually sharper and often more informative for downstream tasks such as classification or change detection~\cite{Li2022SRGANRemoteSensing}.

Furthermore, Super-Resolution GANs (SRGANs) have been applied to upsample low-resolution satellite imagery, enhancing the visibility of small-scale features~\cite{Ledig2017SRGAN,Li2022SRGANRemoteSensing}.
In SAR and RS contexts, such models can improve the apparent spatial resolution while preserving, or even accentuating, structural features like linear infrastructures, small water bodies, or fine topographic details~\cite{Donini2024SuperResRS}.
Recent work by Donini et al. (FBK, UniTn) specifically addresses super-resolution of radargrams, demonstrating that GANs can learn to recover fine stratigraphic details and reduce the impact of limited range resolution in planetary RS data~\cite{Donini2024SuperResRS}.

\subsection{Sim-to-Real Transfer}
A critical limitation in training DL models for planetary science is the scarcity of labeled \enquote{ground truth} data.
Unlike terrestrial applications where ground validation is feasible, planetary RS observations cannot be easily verified through in-situ measurements such as drill cores or geophysical surveys.
Even expert manual annotations remain subjective, inconsistent across interpreters, and fundamentally limited by the ambiguity inherent to radargram interpretation.
This paucity of reliable labels has motivated the development of \textbf{Simulation-to-Real (Sim-to-Real)} transfer learning approaches~\cite{Zhu_2017,yebasse2025advances}.

The workflow typically involves:
\begin{enumerate}
    \item Generating extensive synthetic datasets using physics-based simulators (e.g., FDTD Maxwell solvers or geometric optics clutter models) that provide complete control over subsurface geometry, dielectric properties, and acquisition parameters~\cite{Warren_2016}.
    \item Training models on this clean, fully labeled synthetic data, which includes ground truth for subsurface interfaces, clutter masks, and noise-free conditions.
    \item Adapting the model to solve tasks in the real domain.
\end{enumerate}

CycleGAN~\cite{Zhu_2017} and related architectures are particularly effective for this adaptation, as they learn to translate between \enquote{Clean Simulation} and \enquote{Noisy Real} domains without requiring perfectly aligned training pairs.
This enables \textbf{controlled data augmentation}: starting from a single simulated radargram, the network can generate multiple realistic variants that preserve the subsurface structure while introducing the full spectrum of real-world clutter statistics, speckle textures, and instrumental effects observed in SHARAD data~\cite{yebasse2025advances}.
The resulting augmented dataset bridges the sim-to-real domain gap, substantially improving model generalization to operational planetary RS observations.


